\bookchapter{Bab 17: Koneksi ponsel dengan perangkat lain}

\chapterintro{1. Kolaborasi dengan komputer: Ponsel dan komputer bukan lagi dua dunia}

\booksection{Sejujurnya, dulu yang paling saya benci adalah memotret di ponsel, ingin menaruhnya di komputer untuk diedit, masih harus mencari kabel data, menunggu transfer. Kemudian saya menemukan, sebenarnya tidak perlu kabel data, integrasi yang baik sepuluh kali lebih cepat dari kabel data.}

Ekosistem Apple (AirDrop + Clipboard Universal)

\begin{itemize}
  \item \textbf{: Jika Anda menggunakan iPhone dengan Mac, AirDrop adalah fitur pertama yang harus Anda pelajari. Di ponsel klik Bagikan AirDrop pilih Mac Anda, foto langsung terkirim, tanpa kuota tanpa Wi-Fi, menggunakan Bluetooth dan Wi-Fi Direct. Ada juga fitur yang mudah diabaikan bernama "Clipboard Universal" — menyalin teks di ponsel, di Mac langsung tempel keluar, seperti sihir.}Android + Windows (Link to Windows)
  \item \textbf{: Microsoft punya aplikasi bernama "Link to Windows", ponsel Samsung sudah built-in, ponsel Android lain pergi ke toko aplikasi untuk mengunduh. Setelah dipasangkan Anda bisa langsung melihat notifikasi ponsel, membalas SMS, menelusuri foto, bahkan menerima panggilan di komputer. Saya sendiri menggunakannya selama dua tahun, saat meeting tidak perlu terus menunduk melihat ponsel, notifikasi muncul di layar komputer sudah cukup.}Kolaborasi multi-layar merek
  \item \textbf{: "Multi-screen Collaboration" Huawei dan "Xiaomi Share" Xiaomi lebih hebat lagi — langsung memproyeksikan layar ponsel ke komputer, Anda bisa menggunakan mouse untuk menyeret file dari ponsel ke komputer, juga bisa men-scroll TikTok, membalas WeChat di komputer. Salin-tempel lintas perangkat, seret file semuanya bisa, pengalaman sangat mulus.}2. Transmisi layar ke TV: Tonton bersama keluarga
\end{itemize}

\booksection{Saat pulang untuk tahun baru Imlek, sepupu perempuan saya ingin memutar video perjalanan wisuda di TV. Setelah berjuang setengah mati, akhirnya menemukan bahwa sebenarnya sangat sederhana:}

iPhone TV

\begin{itemize}
  \item \textbf{: Tarik ke bawah Pusat Kontrol dari kanan atas, klik ikon dua kotak bertumpuk (Screen Mirroring), temukan nama TV Anda, sekali klik terhubung. Prasyaratnya TV dan ponsel terhubung Wi-Fi yang sama, dan TV mendukung AirPlay (sebagian besar TV pintar mendukung).}Android TV
  \item \textbf{: Setiap merek namanya berbeda, Huawei menyebutnya "Transmisi Nirkabel", Xiaomi "Proyeksi", OPPO/vivo "Interaksi Multi-Layar". Umumnya bisa ditemukan di panel tombol cepat, jika tidak ada cari "transmisi layar" di pengaturan.}Transmisi langsung dari aplikasi video
  \item \textbf{: Jika Anda hanya ingin menonton iQiyi, Bilibili, YouTube di TV, sebenarnya lebih sederhana — di antarmuka pemutaran cari ikon kecil persegi dengan simbol Wi-Fi, klik, pilih TV, video langsung berpindah. Saat itu ponsel bisa melakukan hal lain, video terus diputar di TV, tidak saling mempengaruhi.}Bagaimana jika transmisi layar gagal
  \item \textbf{: 90\% kasus adalah tidak terhubung Wi-Fi yang sama. 10\% lainnya adalah fitur transmisi layar TV tidak diaktifkan, cari "Interaksi Multi-Layar" atau "AirPlay" di pengaturan TV.}3. Dengan jam tangan pintar/gelang kebugaran: Manajer kesehatan di pergelangan tangan
\end{itemize}

\booksection{Saya membelikan ibu saya gelang kebugaran, awalnya dia berkata “Saya tidak olahraga beli ini untuk apa”. Setelah digunakan dua minggu, dia menemukan bahwa dia hanya tidur lima jam lebih setiap hari, mulai secara aktif menyesuaikan jadwal tidur. Ini adalah nilai paling praktis dari jam tangan/gelang kebugaran — mengubah data kesehatan yang tidak terlihat menjadi angka yang terlihat.}

Filter notifikasi

\begin{itemize}
  \item \textbf{: Saat ponsel bergetar terus di tas atau saku, angkat pergelangan tangan untuk melihat — apakah WeChat dari bos atau promosi Pinduoduo, tentukan dalam satu detik, pesan yang tidak perlu mengeluarkan ponsel biarkan saja.}Sinkronisasi data kesehatan
  \item \textbf{: Jam tangan memantau detak jantung, oksigen darah, tidur, langkah 24 jam, data ini akan otomatis dijumlahkan ke aplikasi Kesehatan di ponsel. Seiring waktu Anda dapat melihat tren Anda — apakah bulan ini berjalan terlalu sedikit, apakah kualitas tidur minggu ini membaik.}Sinkronisasi kartu transportasi/kartu akses
  \item \textbf{: Kartu bus dan kartu akses yang terikat di ponsel, sebagian besar merek akan otomatis disinkronkan ke jam tangan. Saat naik metro tidak perlu mengeluarkan ponsel, pergelangan tangan ditempelkan ke gerbang langsung lewat, benar-benar ada perasaan "masa depan telah tiba".}Pencari ponsel ajaib
  \item \textbf{: Jam tangan umumnya memiliki fungsi "temukan ponsel", klik sekali ponsel akan berbunyi. Saya sendiri menggunakan fungsi ini jauh melampaui ekspektasi — ponsel diletakkan sembarangan lalu lupa di mana.}4. Dengan rumah pintar: Ponsel adalah "saklar utama" rumah
\end{itemize}

\booksection{Rumah pintar terdengar rumit, tetapi sebenarnya hanya tiga hal: kontrol jarak jauh dengan ponsel, atur aturan otomatisasi, kontrol suara.}

Mari saya ceritakan beberapa penggunaan praktis saya sendiri, Anda akan tahu betapa praktisnya hal ini:

AC

\begin{itemize}
  \item \textbf{: Saat perjalanan pulang di musim panas, nyalakan AC dari jarak jauh dengan ponsel sebelum tiba, masuk pintu langsung dingin. Beli colokan pendamping AC seratus yuan lebih bisa mewujudkannya, tidak perlu ganti AC.}Lampu
  \item \textbf{: Sebelum tidur berbaring di tempat tidur, malas bangun untuk mematikan lampu, klik di ponsel atau ucapkan "Hey Siri/Xiao Ai, matikan lampu".}Robot penyedot debu
  \item \textbf{: Sebelum keluar klik "mulai pembersihan" di ponsel, pulang lantai sudah bersih. Juga bisa diatur pembersihan terjadwal setiap hari.}Kamera
  \item \textbf{: Jika di rumah ada orang tua, anak atau hewan peliharaan, di ponsel bisa melihat kondisi rumah kapan saja.}Pilih ekosistem mana? Satu kalimat ringkasan:
\end{itemize}

Gunakan iPhone prioritaskan perangkat yang mendukung Apple HomeKit (sedikit lebih mahal tetapi perlindungan privasi baik).

\begin{itemize}
  \item Banyak perangkat mencari nilai untuk uang ekosistem Xiaomi Mi Home paling lengkap, lampu, colokan, penyedot debu semuanya ada.
  \item Gunakan ponsel Huawei prioritaskan HarmonyOS Connect, pengalaman terbaik dengan perangkat Huawei.
  \item 5. Earphone Bluetooth dan speaker: Kebebasan nirkabel
\end{itemize}

\booksection{Beralih dari earphone berkabel ke earphone Bluetooth, bagi saya adalah peningkatan pengalaman ponsel yang paling jelas. Tidak lagi mengeluarkan gulungan kabel berantakan dari tas, juga tidak perlu khawatir kabel tersangkut di pegangan pintu.}

Pasangan sangat sederhana

\begin{itemize}
  \item \textbf{: Earphone saat pertama dinyalakan umumnya otomatis masuk mode pasangan (lampu berkedip), ponsel buka Bluetooth bisa mencarinya, klik sekali terhubung. Setelah itu setiap kali dikeluarkan otomatis terhubung, tidak perlu melakukan apa-apa.}Merek yang sama ada bonus
  \item \textbf{: AirPods dengan iPhone akan otomatis popup menampilkan daya baterai, juga bisa beralih mulus antara beberapa perangkat Apple. Huawei FreeBuds dengan ponsel Huawei, earphone Xiaomi dengan ponsel Xiaomi juga memiliki popup dan optimasi serupa.}Speaker Bluetooth
  \item \textbf{: Saat kumpul keluarga, berkemah outdoor, sebuah speaker Bluetooth bisa melipatgandakan suasana. Setelah terhubung, ponsel adalah stasiun DJ Anda.}Perhatikan latensi
  \item \textbf{: Saat bermain game jika menggunakan earphone Bluetooth, mungkin ada latensi ringan (tembakan setengah detik baru terdengar suara tembakan). Jika sensitif terhadap latensi, gunakan earphone berkabel, atau beli earphone Bluetooth yang mendukung "mode latensi rendah".}6. Dengan sistem mobil: Mitra aman saat berkendara
\end{itemize}

\begin{tcbinfo}{Teknologi UWB Ultra-Wideband — "Sense of Space" yang Lebih Akurat dari Bluetooth}
Selain Bluetooth, di ponsel baru juga tersembunyi chip UWB (Ultra-Wideband). Kemampuannya adalah\textbf{pengukuran jarak presisi dan posisi spasial}— bisa menghitung seberapa jauh perangkat lain dari Anda, di arah mana, akurasi sampai tingkat sentimeter. Bluetooth hanya bisa memberi tahu “di dekat”, UWB bisa memberi tahu “di 1.2 meter depan kiri Anda”. Aplikasi umum:

\begin{itemize}
  \item \textbf{AirTag/Find Presisi}: Saat iPhone mencari AirTag menampilkan navigasi panah "7 meter jauhnya, jalan ke kiri", mengandalkan UWB. Bluetooth biasa hanya bisa memberi tahu "dalam jangkauan atau tidak".
  \item \textbf{Digital Car Key}: iPhone dan flagship Android yang dilengkapi UWB, bisa menjadi car key mendekati mobil auto unlock, menjauh auto lock, juga bisa membedakan posisi pengemudi dan penumpang, anti relay attack.
  \item \textbf{Navigasi Indoor}: Di mall, bandara via beacon UWB mengimplementasikan navigasi "tingkat lorong", jauh lebih akurat dari GPS.
  \item \textbf{Model yang Mendukung}: iPhone 11 dan model setelahnya semua punya chip UWB (U1/U2); kubu Android Xiaomi, Samsung, vivo, OPPO sebagian flagship juga sudah dilengkapi.
\end{itemize}

\end{tcbinfo}

\begin{tcbinfo}{LE Audio dan Auracast — Standar Generasi Baru Audio Bluetooth}
Audio Bluetooth klasik (A2DP) yang dipakai dua puluh tahun sedang digantikan oleh\textbf{LE Audio}(Low Energy Audio), ini standar audio generasi baru sejak Bluetooth 5.2:

\begin{itemize}
  \item \textbf{Lebih hemat daya, kualitas suara lebih baik}: Memakai codec LC3 yang baru, dibandingkan SBC standar lama kualitas suara lebih baik, konsumsi daya lebih rendah, battery life earphone meningkat nyata.
  \item \textbf{Dual Path Audio}: Earbud kiri dan kanan transmisi independen, switch mode mono tidak terputus.
  \item \textbf{Audio Broadcast Auracast}: Fitur andalan LE Audio. Satu ponsel/TV bisa sekaligus ke\textbf{tak terhingga}earphone broadcast audio — ruang tunggu bandara, museum, gym, kelas semua bisa pakai satu sumber transmit, banyak orang dengan earphone sendiri mendengarkan source yang sama.
  \item \textbf{Status saat ini}: 2024-2026 earphone dan ponsel flagship yang baru diluncurkan umumnya mendukung LE Audio, perangkat lama butuh upgrade hardware tidak bisa di-unlock software. Saat beli earphone baru bisa perhatikan apakah ada label LE Audio/Auracast.
\end{itemize}

\end{tcbinfo}

\booksection{Melihat ponsel saat berkendara adalah ilegal, juga mematikan. Tetapi terkadang benar-benar perlu navigasi, menerima panggilan, memutar musik. Saat itulah CarPlay (Apple) dan Android Auto (Android) berguna.}

Cara menghubungkan

\begin{itemize}
  \item \textbf{: Jika mobil Anda mendukung (sebagian besar mobil baru beberapa tahun terakhir mendukung), gunakan kabel data untuk mencolokkan ponsel ke port USB mobil, layar mobil akan otomatis beralih ke antarmuka CarPlay atau Android Auto. Beberapa mobil baru mendukung koneksi nirkabel, lebih praktis.}Apa yang bisa dilakukan
  \item \textbf{: Navigasi (jauh lebih baik dari navigasi mobil), menelepon (mendengarkan dari setir), mendengarkan musik (aplikasi musik di ponsel langsung diproyeksikan ke mobil), mengirim WeChat dengan suara (Siri atau Xiao Ai membantu membaca pesan dan membalas).}Mobil lama juga bisa
  \item \textbf{: Jika mobil Anda lebih lama tidak mendukung, beli layar CarPlay eksternal beberapa ratus yuan, tempel di dashboard, pengalamannya hampir sama.}7. Keamanan dan privasi dalam integrasi
\end{itemize}

\booksection{Semakin banyak perangkat terintegrasi, semakin besar permukaan risiko kebocoran privasi. Beberapa hal yang perlu diperhatikan:}

Sebelum transmisi layar di tempat umum, periksa dulu apakah ada pesan pribadi di bilah notifikasi ponsel, jangan biarkan seluruh keluarga melihat hal yang tidak seharusnya dilihat.

\begin{itemize}
  \item Sebelum memberikan atau menjual perangkat ke orang lain, lepaskan ikatan dulu — temukan perangkat tersebut di pengaturan ponsel, pilih "Hapus" atau "Batalkan Pasangan".
  \item Kamera rumah pintar, beli yang bermerek, terenkripsi, jangan beli yang murah puluhan yuan tanpa merek. Terlalu banyak berita kamera di-hack untuk live streaming.
  \item Mari coba: Proyeksikan layar ponsel ke TV, dan atur satu skenario pintar
\end{itemize}

\begin{tcbpractice}{Pastikan ponsel dan TV terhubung Wi-Fi yang sama.}
\begin{enumerate}
  \item Di TV temukan fungsi "Screen Mirroring" atau "AirPlay" dan aktifkan.
  \item Di ponsel: iPhone dari Pusat Kontrol klik "Screen Mirroring"; Android tarik ke bawah cari "Transmisi Nirkabel".
  \item Pilih foto favorit Anda proyeksikan ke TV, biarkan keluarga melihat.
  \item Jika ada perangkat rumah pintar, coba buat "Mode Keluar Rumah" di aplikasi ponsel: satu klik matikan lampu, matikan AC, mulai robot penyedot debu.
  \item Bab 18: Pencadangan dan migrasi data
\end{enumerate}

\end{tcbpractice}

\begin{tcbinfo}{Studi Kasus 1: Praktik Optimasi Battery Life}
\begin{itemize}
  \item \textbf{Skenario}: Pengguna Xiaomi 14 battery life buruk, setengah hari sudah habis.
  \item \textbf{Langkah}: (1) Settings Battery Lihat ranking konsumsi daya (2) Identifikasi App background abnormal (3) Matikan 5G switch ke 4G (4) Batasi background (5) Aktifkan power saving mode (6) Saat charge lepas case untuk pendinginan.
  \item \textbf{Hasil yang Diharapkan}: Battery life memanjang 30-50\%, data battery health bisa dibaca.
\end{itemize}

\end{tcbinfo}
